\Chapter{36}

\Verse{1}
{
    \zA{话说贾母自王夫人处回来，见宝玉一日好似一日，心中自是欢喜。}
}
{
    \eA{Ever since dowager lady Chia's return from Madame Wang's quarters, for we
        will now take up the string of our narrative, she naturally felt happier in
        her mind as she saw that Pao-yü improved from day to day;}
}
{
}

\Verse{2}
{
    \zA{因怕将来贾 政又叫他，遂命人将贾政的亲随小厮头儿唤来，吩咐：“以后倘有会人待客诸样的
        事，你老爷要叫宝玉，你不用上来传话，就回他说我说的：一则打重了，得着实将 养几个月才走得；二则他的星宿不利，祭了星，不见外人，过了八月，才许出二门。”}
}
{
    \eA{but nervous lest Chia Cheng should again in the future send for him, she
        lost no time in bidding a servant summon a head-page, a constant attendant
        upon Chia Cheng, to come to her, and in impressing upon him various orders.
        "Should," she enjoined him, "anything turn up henceforward connected with
        meeting guests,}
}
{
}

\Verse{3}
{
    \zA{那小厮头儿听了领命而去。}
}
{
    \eA{entertaining visitors and other such matters, and your master mean to send
        for Pao-yü,}
}
{
}

\Verse{4}
{
    \zA{贾母又命李嬷嬷袭人等来将此话说与宝玉，使他放心。}
}
{
    \eA{you can dispense with going to deliver the message. Just you tell him that I
        say that after the severe thrashing he has had, great care must be first
        taken of him during several months before he can be allowed to walk;}
}
{
}

\Verse{5}
{
    \zA{那宝玉素日本就懒与士大夫诸男人接谈，又最厌峨冠礼服贺吊往还等事，今日得了 这句话，越发得意了，不但将亲戚朋友一概杜绝了，而且连家庭中晨昏定省一发都
        随他的便了。}
}
{
    \eA{and that, secondly, his constellation is unpropitious and that he could not
        see any outsider, while sacrifices are being offered to the stars; that I
        won't have him therefore put his foot beyond the second gate before the
        expiry of the eighth moon." The head-page listened patiently to her
        instructions, and, assenting to all she had to say, he took his leave. Old
        lady Chia thereupon also sent for nurse Li, Hsi Jen and the other
        waiting-maids and recommended them to tell Pao-yü about her injunctions so
        that he might be able to quiet his mind.}
}
{
}

\Verse{6}
{
    \zA{日日只在园中游玩坐卧，不过每日一清早到贾母王夫人处走走就回来 了，却每日甘心为诸丫头充役，倒也得十分消闲日月。}
}
{
    \eA{Pao-yü had always had a repugnance for entertaining high officials and men
        in general, and the greatest horror of going in official hat and ceremonial
        dress, to offer congratulations, or express condolences, to pay calls,
        return visits, or perform other similar conventionalities, but upon receipt
        on the present occasion of this message,}
}
{
}

\Verse{7}
{
    \zA{或如宝钗辈有时见机劝导， 反生起气来，只说：“好好的一个清净洁白女子，也学的钓名沽誉，入了国贼禄鬼 之流。}
}
{
    \eA{he became so much the more confirmed in his dislikes that not only did he
        suspend all intercourse with every single relative and friend, but even went
        so far as to study more than he had ever done before, his own caprices in
        the fulfilment of those morning and evening salutations due to the senior
        members of his family. Day after day he spent in the garden,}
}
{
}

\Verse{8}
{
    \zA{这总是前人无故生事，立意造言，原为引导后世的须眉浊物。}
}
{
    \eA{doing nothing else than loafing about, sitting down here, or reclining
        there. Of a morning, he would, as soon as it was day, stroll as far as the
        quarters of dowager lady Chia and Madame Wang,}
}
{
}

\Verse{9}
{
    \zA{不想我生不幸， 亦且琼闺绣阁中亦染此风，真真有负天地钟灵毓秀之德了！”}
}
{
    \eA{to repair back, however, in no time. Yet ever ready was he every day that
        went by to perform menial services for any of the waiting-maids. He, in
        fact, wasted away in the most complete \_dolce far niente\_ days as well as
        months.}
}
{
}

\Verse{10}
{
    \zA{众人见他如此，也都 不向他说正经话了。}
}
{
    \eA{If perchance Pao-ch'ai or any other girl of the same age as herself found at
        any time an opportunity to give him advice,}
}
{
}

\Verse{11}
{
    \zA{独有黛玉自幼儿不曾劝他去立身扬名，所以深敬黛玉。}
}
{
    \eA{he would, instead of taking it in good part, fly into a huff. "A pure and
        spotless maiden," he would say, "has likewise gone and deliberately imitated
        those persons,}
}
{
}

\Verse{12}
{
    \zA{闲言少述。}
}
{
    \eA{whose aim is to fish for reputation and to seek praise;}
}
{
}

\Verse{13}
{
    \zA{如今且说凤姐自见金钏儿死后，忽见几家仆人常来孝敬他些东西， 又不时的来请安奉承，自己倒生了疑惑，不知何意。}
}
{
    \eA{that set of government thieves and salaried devils. This result entirely
        arises from the fact that there have been people in former times, who have
        uselessly stirred up trouble and purposely fabricated stories with the
        primary object of enticing the filthy male creatures, who would spring up in
        future ages, to follow in their steps!}
}
{
}

\Verse{14}
{
    \zA{这日又见人来孝敬他东西，因 晚间无人时笑问平儿。}
}
{
    \eA{And who would have thought it, I have had the misfortune of being born a
        masculine being! But, even those beautiful girls, in the female apartments,}
}
{
}

\Verse{15}
{
    \zA{平儿冷笑道：“奶奶连这个都想不起来了？}
}
{
    \eA{have been so contaminated by this practice that verily they show themselves
        ungrateful for the virtue of Heaven and Earth, in endowing them with
        perception, and in rearing them with so much comeliness."}
}
{
}

\Verse{16}
{
    \zA{我猜他们的女孩 儿都必是太太屋里的丫头，如今太太屋里有四个大的，一个月一两银子的分例，下 剩的都是一个月只几百钱。}
}
{
    \eA{Seeing therefore what an insane mania possessed him, not one of his cousins
        came forward to tender him one proper word of counsel. Lin Tai-yü was the
        only one of them, who, from his very infancy, had never once admonished him
        to strive and make a position and attain fame, so thus it was that he
        entertained for Tai-yü profound consideration.}
}
{
}

\Verse{17}
{
    \zA{如今金钏儿死了，必定他们要弄这一两银子的窝儿呢。”}
}
{
    \eA{But enough of minor details. We will now turn our attention to lady Feng.
        Soon after the news of Chin Ch'uan-erh's death reached her,}
}
{
}

\Verse{18}
{
    \zA{凤姐听了，笑道：“是了，是了，倒是你想的不错。}
}
{
    \eA{she saw that domestics from various branches of the family paid her frequent
        visits at most unexpected hours, and presented her a lot of things,}
}
{
}

\Verse{19}
{
    \zA{只是这起人也太不知足。}
}
{
    \eA{and that they courted her presence at most unseasonable moments,}
}
{
}

\Verse{20}
{
    \zA{钱也 赚够了，苦事情又摊不着他们，弄个丫头搪塞身子儿也就罢了，又要想这个巧宗儿! 他们几家的钱也不是容易花到我跟前的，这可是他们自寻。}
}
{
    \eA{to pay their compliments and adulate her, and she begun to harbour
        suspicions, in her own mind, as she little knew what their object could
        possibly be. On this date, she again noticed that some of them had brought
        their gifts, so, when evening arrived, and no one was present, she felt
        compelled to inquire jocosely of P'ing Erh what their aim could be.}
}
{
}

\Verse{21}
{
    \zA{送什么我就收什么，横 竖我有主意。”}
}
{
    \eA{"Can't your ladyship fathom even this?" P'ing Erh answered with a sardonic
        smile. "Why, their daughters must,}
}
{
}

\Verse{22}
{
    \zA{凤姐儿安下这个心，所以只管耽延着，等那些人把东西送足了，然 后乘空方回王夫人。}
}
{
    \eA{I fancy, be servant-girls in Madame Wang's apartments! For her ladyship's
        rooms four elderly girls are at present allotted with a monthly allowance of
        one tael; the rest simply receiving several hundreds of cash each month; so
        now that Chin Ch'uan-erh is dead and gone,}
}
{
}

\Verse{23}
{
    \zA{这日午间，薛姨妈、宝钗、黛玉等正在王夫人屋里，大家吃西瓜。}
}
{
    \eA{these people must, of course, be anxious to try their tricks and get this
        one-tael job!" Hearing this, lady Feng smiled a significant smile. "That's
        it. Yes, that's it!"}
}
{
}

\Verse{24}
{
    \zA{凤姐儿得便 回王夫人道：“自从玉钏儿的姐姐死了，太太跟前少着一个人，太太或看准了那个 丫头，就吩咐了，下月好发放月钱。”}
}
{
    \eA{she exclaimed. "You've really suggested the idea to my mind! From all
        appearances, these people are a most insatiable lot; for they make quite
        enough in the way of money! And as for any business that requires a little
        exertion, why they are never ready to bear a share of it! They make use of
        their girls as so many tools to shove their own duties upon.}
}
{
}

\Verse{25}
{
    \zA{王夫人听了，想了一想道：“依我说，什么 是例，必定四个五个的？}
}
{
    \eA{Yet one overlooks that. But must they too have designs upon this job? Never
        mind! These people cannot easily afford to spend upon me the money they do.
        But they bring this upon their own selves,}
}
{
}

\Verse{26}
{
    \zA{够使就罢了。}
}
{
    \eA{so I'll keep every bit of thing they send.}
}
{
}

\Verse{27}
{
    \zA{竟可以免了罢。”}
}
{
    \eA{I've, after all, resolved how to act in the matter!"}
}
{
}

\Verse{28}
{
    \zA{凤姐笑道：“论理，太太 说的也是；只是原是旧例。}
}
{
    \eA{Having arrived at this decision, lady Feng purely and simply protracted the
        delay until all the women had sent her enough to satisfy her,}
}
{
}

\Verse{29}
{
    \zA{别人屋里还有两个呢，太太倒不按例了。}
}
{
    \eA{when she at last suited her own convenience and spoke to Madame Wang (on the
        subject of the vacant post).}
}
{
}

\Verse{30}
{
    \zA{况且省下一两 银子，也有限的。”}
}
{
    \eA{Mrs. Hsüeh and her daughter were sitting one day, at noon, in Madame Wang's
        quarters,}
}
{
}

\Verse{31}
{
    \zA{王夫人听了，又想了想道：“也罢，这个分例只管关了来，不 用补人，就把这一两银子给他妹妹玉钏儿罢。}
}
{
    \eA{together with Lin Tai-yü and the other girls, when lady Feng found an
        opportunity and broached the topic with Madame Wang. "Ever since," she said,
        "sister Chin Ch'uan-erh's death, there has been one servant less in your
        ladyship's service. But you may possibly have set your choice upon some
        girl;}
}
{
}

\Verse{32}
{
    \zA{他姐姐伏侍了我一场，没个好结果， 剩下他妹妹跟着我，吃个双分儿也不为过。”}
}
{
    \eA{if so, do let me know who it is, so that I may be able to pay her her
        monthly wages." This reminder made Madame Wang commune with her own self. "I
        fancy," she remarked; "that the custom is that there should be four or five
        of them;}
}
{
}

\Verse{33}
{
    \zA{凤姐答应着，回头望着玉钏儿笑道： “大喜，大喜！”}
}
{
    \eA{but as long as there are enough to wait upon me, I don't mind, so we can
        really dispense with another." "What you say is, properly speaking,}
}
{
}

\Verse{34}
{
    \zA{玉钏儿过来磕了头。}
}
{
    \eA{perfectly correct," smiled lady Feng; "but it's an old established custom.}
}
{
}

\Verse{35}
{
    \zA{王夫人又问道：“正要问你：如今赵姨娘周姨娘的月例多少？”}
}
{
    \eA{There are still a couple to be found in other people's rooms and won't you,
        Madame, conform with the rule? Besides, the saving of a tael is a small
        matter." After this argument, Madame Wang indulged in further thought.}
}
{
}

\Verse{36}
{
    \zA{凤姐道：“那 是定例，每人二两。}
}
{
    \eA{"Never mind," she then observed, "just you bring over this allowance and pay
        it to me.}
}
{
}

\Verse{37}
{
    \zA{赵姨娘有环兄弟的二两，共是四两，另外四串钱。”}
}
{
    \eA{And there will be no need to supply another girl. I'll hand over this tael
        to her younger sister, Yü Ch'uan-erh, and finish with it. Her elder sister
        came to an unpleasant end,}
}
{
}

\Verse{38}
{
    \zA{王夫人道： “月月可都按数给他们？”}
}
{
    \eA{after a long term of service with me; so if the younger sister, she leaves
        behind in my employ, receives a double share,}
}
{
}

\Verse{39}
{
    \zA{凤姐见问得奇，忙道：“怎么不按数给呢！”}
}
{
    \eA{it won't be any too excessive." Lady Feng expressed her approval and turning
        round she said smilingly to Yü Ch'uan-erh: "I congratulate you, I
        congratulate you!"}
}
{
}

\Verse{40}
{
    \zA{王夫人道： “前儿恍惚听见有人抱怨，说短了一串钱，什么原故？”}
}
{
    \eA{Yü Ch'uan-erh thereupon crossed over and prostrated herself. "I just want to
        ask you," Madame Wang went on to inquire, "how much Mrs. Chao and Mrs. Chou
        are allowed monthly?" "They have a fixed allowance,"}
}
{
}

\Verse{41}
{
    \zA{凤姐忙笑道：“姨娘们的 丫头月例，原是人各一吊钱，从旧年他们外头商量的，姨娘们每位丫头，分例减半， 人各五百钱。}
}
{
    \eA{answered lady Feng, "each of them draws two taels. But Mrs. Chao gets two
        taels for cousin Chia Huan, so hers amounts in all to four taels; besides
        these, four strings of cash." "Are they paid in full month after month?"
        Madame Wang inquired. Lady Feng thought the question so very strange that
        she hastened to exclaim by way of reply:}
}
{
}

\Verse{42}
{
    \zA{每位两个丫头，所以短了一吊钱。}
}
{
    \eA{"How are they not paid in full?" "The other day," Madame Wang proceeded, "I
        heard a faint rumour that there was some one,}
}
{
}

\Verse{43}
{
    \zA{这事其实不在我手里，我倒乐得给 他们呢，只是外头扣着，这里我不过是接手儿，怎么来怎么去，由不得我做主。}
}
{
    \eA{who complained in an aggrieved way that she had got a string short. How and
        why is this?" "The monthly allowances of the servant-girls, attached to the
        secondary wives," lady Feng hurriedly added with a smile, "amounted
        originally to a tiao each,}
}
{
}

\Verse{44}
{
    \zA{我 倒说了两三回，仍旧添上这两分儿为是，他们说了‘只有这个数儿’，叫我也难再 说了。}
}
{
    \eA{but ever since last year, it was decided, by those people outside, that the
        shares of each of those ladies' girls should be reduced by half, that is,
        each to five hundred cash; and, as each lady has a couple of servant-girls,}
}
{
}

\Verse{45}
{
    \zA{如今我手里给他们，每月连日子都不错。}
}
{
    \eA{they receive therefore a tiao short. But for this, they can't bear me a
        grudge. As far as I'm concerned,}
}
{
}

\Verse{46}
{
    \zA{先时候儿在外头关，那个月不打饥 荒，何曾顺顺溜溜的得过一遭儿呢。”}
}
{
    \eA{I would only be too glad to let them have it; but our people outside will
        again disallow it; so is it likely that I can authorise any increase, pray?
        In this matter of payments I merely receive the money,}
}
{
}

\Verse{47}
{
    \zA{王夫人听说，就停了半晌，又问：“老太太 屋里几个一两的？”}
}
{
    \eA{and I've nothing to do with how it comes and how it goes. I nevertheless
        recommended, on two or three occasions, that it would be better if these two
        shares were again raised to the old amount;}
}
{
}

\Verse{48}
{
    \zA{凤姐道：“八个。}
}
{
    \eA{but they said that there's only that much money,}
}
{
}

\Verse{49}
{
    \zA{如今只有七个，那一个是袭人。”}
}
{
    \eA{so that I can't very well volunteer any further suggestions! Now that the
        funds are paid into my hands,}
}
{
}

\Verse{50}
{
    \zA{王夫人说： “这就是了。}
}
{
    \eA{I give them to them every month, without any irregularity of even so much as
        a day.}
}
{
}

\Verse{51}
{
    \zA{你宝兄弟也并没有一两的丫头，袭人还算老太太房里的人。”}
}
{
    \eA{When payments hitherto were effected outside, what month were they not short
        of money? And did they ever, on any single instance, obtain their pay at the
        proper time and date?"}
}
{
}

\Verse{52}
{
    \zA{凤姐笑 道：“袭人还是老太太的人，不过给了宝兄弟使，他这一两银子还在老太太的丫头 分例上领。}
}
{
    \eA{Having heard this explanation, Madame Wang kept silent for a while. Next,
        she proceeded to ask, how many girls there were with dowager lady Chia
        drawing one tael. "Eight of them," rejoined lady Feng, "but there are at
        present only seven; the other one is Hsi Jen."}
}
{
}

\Verse{53}
{
    \zA{如今说因为袭人是宝玉的人，裁了这一两银子，断乎使不得。}
}
{
    \eA{"Quite right," assented Madame Wang. "But your cousin Pao-yü hasn't any maid
        at one tael; for Hsi Jen is still a servant belonging to old lady Chia's
        household." "Hsi Jen," lady Feng smiled,}
}
{
}

\Verse{54}
{
    \zA{若说再添 一个人给老太太，这个还可以裁他。}
}
{
    \eA{"is still our dear ancestor's servant; she's only lent to cousin Pao-yü; so
        that she still receives this tael in her capacity of maid to our worthy
        senior.}
}
{
}

\Verse{55}
{
    \zA{若不裁他，须得环兄弟屋里也添上一个，才公 道均匀了。}
}
{
    \eA{Any proposal, therefore, that might now be made, that this tael should, as
        Hsi Jen is Pao-yü's servant, be curtailed, can, on no account,}
}
{
}

\Verse{56}
{
    \zA{就是晴雯、麝月他们七个大丫头，每月人各月钱一吊，佳蕙他们八个小 丫头们，每月人各月钱五百，还是老太太的话，别人也恼不得气不得呀。”}
}
{
    \eA{be entertained. Yet, were it suggested that another servant should be added
        to our senior's staff, then in this way one could reduce the tael she gets.
        But if this be not curtailed, it will be necessary to also add a servant in
        cousin Chia Huan's rooms, in order that there should be a fair
        apportionment. In fact, Ch'ing Wen, She Yüeh and the others, numbering seven
        senior maids,}
}
{
}

\Verse{57}
{
    \zA{薛姨妈笑道：“你们只听凤丫头的嘴，倒像倒了核桃车子似的。}
}
{
    \eA{receive each a tiao a month; and Chiao Hui and the rest of the junior maids,
        eight in all, get each five hundred cash per mensem; and this was
        recommended by our venerable ancestor herself; so how can any one be angry
        and feel displeasure?"}
}
{
}

\Verse{58}
{
    \zA{帐也清楚，理 也公道。”}
}
{
    \eA{"Just listen," laughed Mrs. Hsüeh, "to that girl Feng's mouth!}
}
{
}

\Verse{59}
{
    \zA{凤姐笑道：“姑妈，难道我说错了吗？”}
}
{
    \eA{It rattles and rattles like a cart laden with walnuts, which has turned
        topsy-turvy! Yet, her accounts are, from what one can gather,}
}
{
}

\Verse{60}
{
    \zA{薛姨妈笑道：“说的何尝错， 只是你慢着些儿说不省力些？”}
}
{
    \eA{clear enough, and her arguments full of reason." "Aunt," rejoined lady Feng
        smiling, "was I likely, pray, wrong in what I said?" "Who ever said you were
        wrong?"}
}
{
}

\Verse{61}
{
    \zA{凤姐才要笑，忙又忍住了，听王夫人示下。}
}
{
    \eA{Mrs. Hsüeh smiled. "But were you to talk a little slower, wouldn't it be a
        saving of exertion for you?" Lady Feng was about to laugh, but hastily
        checking herself,}
}
{
}

\Verse{62}
{
    \zA{王夫人 想了半日，向凤姐道：“明儿挑一个丫头送给老太太使唤，补袭人，把袭人的一分 裁了。}
}
{
    \eA{she lent an ear to what Madame Wang might have to tell her. Madame Wang
        indulged in thought for a considerable time. Afterwards, facing lady Feng,
        "You'd better," she said, "select a waiting-maid tomorrow and send her over
        to our worthy senior to fill up Hsi Jen's place.}
}
{
}

\Verse{63}
{
    \zA{把我每月的月例，二十两银子里拿出二两银子一吊钱来，给袭人去。}
}
{
    \eA{Then, discontinue that allowance, which Hsi Jen draws, and keep out of the
        sum of twenty taels, allotted to me monthly, two taels and a tiao, and give
        them to Hsi Jen.}
}
{
}

\Verse{64}
{
    \zA{以后凡 事有赵姨娘周姨娘的，也有袭人的，只是袭人的这一分，都从我的分例上匀出来， 不必动官中的就是了。”}
}
{
    \eA{So henceforward what Mrs. Chao and Mrs. Chou will get, Hsi Jen will likewise
        get, with the only difference that the share granted to Hsi Jen, will be
        entirely apportioned out of my own allowance. Mind, therefore, there will be
        no necessity to touch the public funds!"}
}
{
}

\Verse{65}
{
    \zA{凤姐一一的答应了，笑推薛姨妈道：“姑妈听见了？}
}
{
    \eA{Lady Feng acquiesced to each one of her recommendations, and, pushing Mrs.
        Hsüeh, "Aunt," she inquired, "have you heard her proposal?}
}
{
}

\Verse{66}
{
    \zA{我素 日说的话如何？}
}
{
    \eA{What have I all along maintained? Well, my words have actually come out true
        to-day!"}
}
{
}

\Verse{67}
{
    \zA{今儿果然应了。”}
}
{
    \eA{"This should have been accomplished long ago,"}
}
{
}

\Verse{68}
{
    \zA{薛姨妈道：“早就该这么着。}
}
{
    \eA{Mrs. Hsüeh answered. "For without, of course, making any allusion to her
        looks,}
}
{
}

\Verse{69}
{
    \zA{那孩子模样儿不用 说，只是他那行事儿的大方，见人说话儿的和气，里头带着刚硬要强，倒实在难得 的。”}
}
{
    \eA{her way of doing business is liberal; her speech and her relations with
        people are always prompted by an even temper, while inwardly she has plenty
        of singleness of heart and eagerness to hold her own.}
}
{
}

\Verse{70}
{
    \zA{王夫人含泪说道：“你们那里知道袭人那孩子的好处？}
}
{
    \eA{Indeed, such a girl is not easy to come across!" Madame Wang made every
        effort to conceal her tears. "How could you people ever rightly estimate Hsi
        Jen's qualities?"}
}
{
}

\Verse{71}
{
    \zA{比我的宝玉还强十倍 呢!宝玉果然有造化，能够得他长长远远的伏侍一辈子，也就罢了。”}
}
{
    \eA{she observed. "Why, she's a hundred times better than my own Pao-yü. How
        fortunate, in reality, Pao-yü is! Well would it be if he could have her wait
        upon him for the whole length of his life!" "In that case," lady Feng
        suggested,}
}
{
}

\Verse{72}
{
    \zA{凤姐道：“既 这么样，就开了脸，明放他在屋里不好？”}
}
{
    \eA{"why, have her face shaved at once, and openly place her in his room as a
        secondary wife. Won't this be a good plan?" "This won't do!" Madame Wang
        retorted.}
}
{
}

\Verse{73}
{
    \zA{王夫人道：“这不好：一则年轻；二则 老爷也不许；三则宝玉见袭人是他的丫头，纵有放纵的事，倒能听他的劝，如今做 了跟前人，那袭人该劝的也不敢十分劝了。}
}
{
    \eA{"For first and foremost he's of tender years. In the second place, my
        husband won't countenance any such thing! In the third, so long as Pao-yü
        sees that Hsi Jen is his waiting-maid, he may, in the event of anything
        occurring from his having been allowed to run wild,}
}
{
}

\Verse{74}
{
    \zA{如今且浑着，等再过二三年再说。”}
}
{
    \eA{listen to any good counsel she might give him. But were she now to be made
        his secondary wife,}
}
{
}

\Verse{75}
{
    \zA{说毕，凤姐见无话，便转身出来。}
}
{
    \eA{Hsi Jen would not venture to tender him any extreme advice, even when it's
        necessary to do so.}
}
{
}

\Verse{76}
{
    \zA{刚至廊檐下，只见有几个执事的媳妇子正等 他回事呢，见他出来，都笑道：“奶奶今儿回什么事，说了这半天？}
}
{
    \eA{It's better, therefore, to let things stand as they are for the present, and
        talk about them again, after the lapse of another two or three years." At
        the close of these arguments, lady Feng could not put in a word,}
}
{
}

\Verse{77}
{
    \zA{可别热着罢。”}
}
{
    \eA{by way of reply, to refute them,}
}
{
}

\Verse{78}
{
    \zA{凤姐把袖子挽了几挽，？}
}
{
    \eA{so turning round, she left the room. She had no sooner, however,}
}
{
}

\Verse{79}
{
    \zA{着那角门的门槛子，笑道：“这里过堂风，倒凉快，吹一 吹再走。”}
}
{
    \eA{got under the verandah, than she discerned the wives of a number of butlers,
        waiting for her to report various matters to her. Seeing her issue out of
        the room, they with one consent smiled.}
}
{
}

\Verse{80}
{
    \zA{又告诉众人道：“你们说我回了这半日的话，太太把二百年的事都想起 来问我，难道我不说罢？”}
}
{
    \eA{"What has your ladyship had to lay before Madame Wang," they remarked, "that
        you've been talking away this length of time? Didn't you find it hot work?"
        Lady Feng tucked up her sleeves several times. Then resting her foot on the
        step of the side door, she laughed and rejoined:}
}
{
}

\Verse{81}
{
    \zA{又冷笑道：“我从今以后，倒要干几件刻薄事了。}
}
{
    \eA{"The draft in this passage is so cool, that I'll stop, and let it play on me
        a bit before I go on. You people," she proceeded to tell them,}
}
{
}

\Verse{82}
{
    \zA{抱怨 给太太听，我也不怕!糊涂油蒙了心、烂了舌头、不得好死的下作娼妇们，别做娘 的春梦了!明儿一裹脑子扣的日子还有呢。}
}
{
    \eA{"say that I've been talking to her all this while, but Madame Wang conjured
        up all that has occurred for the last two hundred years and questioned me
        about it; so could I very well not have anything to say in reply? But from
        this day forth," she added with a sarcastic smile, "I shall do several mean
        things,}
}
{
}

\Verse{83}
{
    \zA{如今裁了丫头的钱就抱怨了咱们，也不 想想自己也配使三个丫头！”}
}
{
    \eA{and should even (Mrs. Chao and Mrs. Chou) go, out of any ill-will, and tell
        Madame Wang, I won't know what fear is for such stupid, glib-tongued,
        foul-mouthed creatures as they,}
}
{
}

\Verse{84}
{
    \zA{一面骂，一面方走了，自去挑人回贾母话去，不在话 下。}
}
{
    \eA{who are bound not to see a good end! It isn't for them to indulge in those
        fanciful dreams of becoming primary wives, for there, will come soon a day
        when the whole lump sum of their allowance will be cut off!}
}
{
}

\Verse{85}
{
    \zA{却说薛姨妈等这里吃毕西瓜，又说了一回闲话儿，各自散去。}
}
{
    \eA{They grumble against us for having now reduced the perquisites of the
        servant-maids, but they don't consider whether they deserve to have so many
        as three girls to dance attendance on them!" While heaping abuse on their
        heads,}
}
{
}

\Verse{86}
{
    \zA{宝钗与黛玉回至 园中，宝钗要约着黛玉往藕香榭去，黛玉因说还要洗澡，便各自散了。}
}
{
    \eA{she started homewards, and went all alone in search of some domestic to go
        and deliver a message to old lady Chia. But without any further reference to
        her, we will take up the thread of our narrative with Mrs. Hsüeh, and the
        others along with her.}
}
{
}

\Verse{87}
{
    \zA{宝钗独自行 来，顺路进了怡红院，意欲寻宝玉去说话儿，以解午倦。}
}
{
    \eA{During this interval they finished feasting on melons. After some more
        gossip, each went her own way; and Pao-ch'ai, Tai-yü and the rest of the
        cousins returned into the garden.}
}
{
}

\Verse{88}
{
    \zA{不想步入院中，鸦雀无闻， 一并连两只仙鹤在芭蕉下都睡着了。}
}
{
    \eA{Pao-ch'ai then asked Tai-yü to repair with her to the O Hsiang Arbour. But
        Tai-yü said that she was just going to have her bath, so they parted
        company, and Pao-ch'ai walked back all by herself.}
}
{
}

\Verse{89}
{
    \zA{宝钗便顺着游廊，来至房中。}
}
{
    \eA{On her way, she stepped into the I Hung Yüan, to look up Pao-yü and have a
        friendly hobnob with him,}
}
{
}

\Verse{90}
{
    \zA{只见外间床上横 三竖四，都是丫头们睡觉。}
}
{
    \eA{with the idea of dispelling her mid-day lassitude; but, contrary to her
        expectations, the moment she put her foot into the court,}
}
{
}

\Verse{91}
{
    \zA{转过十锦？}
}
{
    \eA{she did not so much as catch the caw of a crow.}
}
{
}

\Verse{92}
{
    \zA{子，来至宝玉的房内，宝玉在床上睡着了， 袭人坐在身旁，手里做针线，傍边放着一柄白犀麈。}
}
{
    \eA{Even the two storks stood under the banana trees, plunged in sleep.
        Pao-ch'ai proceeded along the covered passage and entered the rooms. Here
        she discovered the servant-girls sleeping soundly on the bed of the outer
        apartment; some lying one way, some another;}
}
{
}

\Verse{93}
{
    \zA{宝钗走近前来，悄悄的笑道：“你也过于小心了。}
}
{
    \eA{so turning round the decorated screen, she wended her steps into Pao-yü's
        chamber. Pao-yü was asleep in bed. Hsi Jen was seated by his side,}
}
{
}

\Verse{94}
{
    \zA{这个屋里还有苍蝇蚊子？}
}
{
    \eA{busy plying her needle. Next to her, lay a yak tail. Pao-ch'ai advanced up
        to her.}
}
{
}

\Verse{95}
{
    \zA{还 拿蝇刷子赶什么？”}
}
{
    \eA{"You're really far too scrupulous," she said smilingly in an undertone.}
}
{
}

\Verse{96}
{
    \zA{袭人不防，猛抬头见是宝钗，忙放针线起身，悄悄笑道：“姑 娘来了，我倒不防，唬了一跳。}
}
{
    \eA{"Are there still flies or mosquitos in here? and why do yet use that
        fly-flap for, to drive what away?" Hsi Jen was quite taken by surprise. But
        hastily raising her head, and realising that it was Pao-ch'ai, she hurriedly
        put down her needlework.}
}
{
}

\Verse{97}
{
    \zA{姑娘不知道：虽然没有苍蝇蚊子，谁知有一种小虫 子，从这纱眼里钻进来，人也看不见。}
}
{
    \eA{"Miss," she whispered with a smile, "you came upon me so unawares that you
        gave me quite a start! You don't know, Miss, that though there be no flies
        or mosquitoes there is, no one would believe it, a kind of small insect,
        which penetrates through the holes of this gauze;}
}
{
}

\Verse{98}
{
    \zA{只睡着了咬一口，就像蚂蚁叮的。”}
}
{
    \eA{it is scarcely to be detected, but when one is asleep, it bites just like
        ants do!" "It isn't to be wondered at,"}
}
{
}

\Verse{99}
{
    \zA{宝钗道： “怨不得，这屋子后头又近水，又都是香花儿，这屋子里头又香，这种虫子都是花 心里长的，闻香就扑。”}
}
{
    \eA{Pao-ch'ai suggested, "for the back of these rooms adjoins the water; the
        whole place is also one mass of fragrant flowers, and the interior of this
        room is, too, full of their aroma. These insects grow mostly in the core of
        flowers, so no sooner do they scent the smell of any than they at once rush
        in."}
}
{
}

\Verse{100}
{
    \zA{说着，一面就瞧他手里的针线。}
}
{
    \eA{Saying this, she cast a look on the needlework she (Hsi Jen) held in her
        hands.}
}
{
}

\Verse{101}
{
    \zA{原来是个白绫红里的兜肚， 上面扎着鸳鸯戏莲的花样，红莲绿叶，五色鸳鸯。}
}
{
    \eA{It consisted, in fact, of a belt of white silk, lined with red, and
        embroidered on the upper part with designs representing mandarin ducks,
        disporting themselves among some lotus. The lotus flowers were red,}
}
{
}

\Verse{102}
{
    \zA{宝钗道：“嗳哟，好鲜亮活计。}
}
{
    \eA{the leaves green, the ducks of variegated colours. "Ai-yah!" ejaculated
        Pao-ch'ai,}
}
{
}

\Verse{103}
{
    \zA{这是谁的，也值的费这么大工夫？”}
}
{
    \eA{"what very beautiful work! For whom is this, that it's worth your while
        wasting so much labour on it?"}
}
{
}

\Verse{104}
{
    \zA{袭人向床上？？}
}
{
    \eA{Hsi Jen pouted her lips towards the bed.}
}
{
}

\Verse{105}
{
    \zA{嘴儿。}
}
{
    \eA{"Does a big strapping fellow like this,"}
}
{
}

\Verse{106}
{
    \zA{宝钗笑道：“这么大了， 还带这个？”}
}
{
    \eA{Pao-ch'ai laughed, "still wear such things?" "He would never wear any
        before,"}
}
{
}

\Verse{107}
{
    \zA{袭人笑道：“他原是不带，所以特特的做的好了，叫他看见，由不得 不带。}
}
{
    \eA{Hsi Jen smiled, "that's why such a nice one was specially worked for him, in
        order that when he was allowed to see it, he should not be able to do
        otherwise than use it.}
}
{
}

\Verse{108}
{
    \zA{如今天热，睡觉都不留神，哄他带上了，就是夜里纵盖不严些儿，也就罢了。}
}
{
    \eA{With the present hot weather, he goes to sleep anyhow, but as he has been
        coaxed to wear it, it doesn't matter if even he doesn't cover himself well
        at night.}
}
{
}

\Verse{109}
{
    \zA{你说这一个就用了工夫，还没看见他身上带的那一个呢！”}
}
{
    \eA{You say that I bestow much labour upon this, but you haven't yet seen the
        one he has on!" "It is a lucky thing," Pao-ch'ai observed,}
}
{
}

\Verse{110}
{
    \zA{宝钗笑道：“也亏你耐 烦。”}
}
{
    \eA{smiling, "that you're gifted with such patience." "I've done so much of it
        to-day,"}
}
{
}

\Verse{111}
{
    \zA{袭人道：“今儿做的工夫大了，脖子低的怪酸的。”}
}
{
    \eA{remarked Hsi Jen, "that my neck is quite sore from bending over it. My dear
        Miss," she then urged with a beaming countenance,}
}
{
}

\Verse{112}
{
    \zA{又笑道：“好姑娘，你 略坐一坐，我出去走走就来。”}
}
{
    \eA{"do sit here a little. I'll go out for a turn. I'll be back shortly." With
        these words, she sallied out of the room.}
}
{
}

\Verse{113}
{
    \zA{说着就走了。}
}
{
    \eA{Pao-ch'ai was intent upon examining the embroidery,}
}
{
}

\Verse{114}
{
    \zA{宝钗只顾看着活计便不留心，一蹲身， 刚刚的也坐在袭人方才坐的那个所在。}
}
{
    \eA{so in her absentmindedness, she, with one bend of her body, settled herself
        on the very same spot, which Hsi Jen had recently occupied.}
}
{
}

\Verse{115}
{
    \zA{因又见那个活计实在可爱，不由的拿起针 来，就替他作。}
}
{
    \eA{But she found, on second scrutiny, the work so really admirable, that
        impulsively picking up the needle, she continued it for her.}
}
{
}

\Verse{116}
{
    \zA{不想黛玉因遇见湘云，约他来与袭人道喜，二人来至院中。}
}
{
    \eA{At quite an unforeseen moment—for Lin Tai-yü had met Shih Hsiang-yün and
        asked her to come along with her and present her congratulations to Hsi
        Jen—these two girls made their appearance in the court. Finding the whole
        place plunged in silence,}
}
{
}

\Verse{117}
{
    \zA{见静悄悄的，湘云 便转身先到厢房里去找袭人去了。}
}
{
    \eA{Hsiang-yün turned round and betook herself first into the side-rooms in
        search of Hsi Jen. Lin Tai-yü, meanwhile, walked up to the window from
        outside,}
}
{
}

\Verse{118}
{
    \zA{那黛玉却来至窗外，隔着窗纱往里一看，只见宝 玉穿着银红纱衫子，随便睡着在床上，宝钗坐在身旁做针线，傍边放着蝇刷子。}
}
{
    \eA{and peeped in through the gauze frame. At a glance, she espied Pao-yü, clad
        in a silvery-red coat, lying carelessly on the bed, and Pao-ch'ai, seated by
        his side, busy at some needlework, with a fly-brush resting by her side.}
}
{
}

\Verse{119}
{
    \zA{黛 玉见了这个景况，早已呆了，连忙把身子一躲，半日又握着嘴笑，却不敢笑出来， 便招手儿叫湘云。}
}
{
    \eA{As soon as Lin Tai-yü became conscious of the situation, she immediately
        slipped out of sight, and stopping her mouth with one hand, as she did not
        venture to laugh aloud, she waved her other hand and beckoned to Hsiang-yün.
        The moment Hsiang-yün saw the way she went on,}
}
{
}

\Verse{120}
{
    \zA{湘云见他这般，只当有什么新闻，忙也来看，才要笑，忽然想起 宝钗素日待他厚道，便忙掩住口。}
}
{
    \eA{she concluded that she must have something new to impart to her, and she
        approached her with all promptitude. At the sight, which opened itself
        before her eyes, she also felt inclined to laugh. Yet the sudden
        recollection of the kindness,}
}
{
}

\Verse{121}
{
    \zA{知道黛玉口里不让人，怕他取笑，便忙拉过他来， 道：“走罢。}
}
{
    \eA{with which Pao-ch'ai had always dealt towards her, induced her to quickly
        seal her lips. And knowing well enough that Tai-yü never spared any one with
        her mouth,}
}
{
}

\Verse{122}
{
    \zA{我想起袭人来，他说晌午要到池子里去洗衣裳，想必去了，咱们找他 去罢。”}
}
{
    \eA{she was seized with such fear lest she should jeer at them, that she
        immediately dragged her past the window. "Come along!" she observed. "Hsi
        Jen, I remember, said that she would be going at noon to wash some clothes
        at the pond.}
}
{
}

\Verse{123}
{
    \zA{黛玉心下明白，冷笑了两声，只得随他走了。}
}
{
    \eA{I presume she's there already so let's go and join her." Tai-yü inwardly
        grasped her meaning, but, after indulging in a couple of sardonic smiles,}
}
{
}

\Verse{124}
{
    \zA{这里宝钗只刚做了两三个花瓣，忽见宝玉在梦中喊骂，说：“和尚道士的话如 何信得？}
}
{
    \eA{she had no alternative but to follow in her footsteps. Pao-ch'ai had, during
        this while, managed to embroider two or three petals, when she heard Pao-yü
        begin to shout abusingly in his dreams.}
}
{
}

\Verse{125}
{
    \zA{什么‘金玉姻缘’？}
}
{
    \eA{"How can," he cried, "one ever believe what bonzes and Taoist priests say?}
}
{
}

\Verse{126}
{
    \zA{我偏说‘木石姻缘’！”}
}
{
    \eA{What about a match between gold and jade?}
}
{
}

\Verse{127}
{
    \zA{宝钗听了这话，不觉怔了。}
}
{
    \eA{My impression is that it's to be a union between a shrub and a stone!"}
}
{
}

\Verse{128}
{
    \zA{忽 见袭人走进来，笑道：“还没醒呢吗？”}
}
{
    \eA{Hsüeh Pao-ch'ai caught every single word uttered by him and fell
        unconsciously in a state of excitement. Of a sudden, however, Hsi Jen
        appeared on the scene.}
}
{
}

\Verse{129}
{
    \zA{宝钗摇头。}
}
{
    \eA{"Hasn't he yet woke up?"}
}
{
}

\Verse{130}
{
    \zA{袭人又笑道：“我才碰见林姑 娘史大姑娘，他们进来了么？”}
}
{
    \eA{she inquired. Pao-ch'ai nodded her head by way of reply. "I just came
        across," Hsi Jen smiled, "Miss Lin and Miss Shih.}
}
{
}

\Verse{131}
{
    \zA{宝钗道：“没见他们进来。”}
}
{
    \eA{Did they happen to come in?" "I didn't see them come in,"}
}
{
}

\Verse{132}
{
    \zA{因向袭人笑道：“他 们没告诉你什么？”}
}
{
    \eA{Pao-ch'ai answered. "Did they tell you anything?" she next smilingly asked
        of Hsi Jen.}
}
{
}

\Verse{133}
{
    \zA{袭人红了脸，笑道：“总不过是他们那些玩话，有什么正经说 的。”}
}
{
    \eA{Hsi Jen blushed and laughed significantly. "They simply came out with some
        of those jokes of theirs," she explained. "What decent things could such as
        they have had to tell me?" "They made insinuations to-day," Pao-ch'ai
        laughed, "which are anything but a joke!}
}
{
}

\Verse{134}
{
    \zA{宝钗笑道：“今儿他们说的可不是玩话，我正要告诉你呢，你又忙忙的出去 了。”}
}
{
    \eA{I was on the point of telling you them, when you rushed away in an awful
        hurry." But no sooner had she concluded, than she perceived a servant, come
        over from lady Feng's part to fetch Hsi Jen. "It must be on account of what
        they hinted,"}
}
{
}

\Verse{135}
{
    \zA{一句话未完，只见凤姐打发人来叫袭人。}
}
{
    \eA{Pao-ch'ai smilingly added. Hsi Jen could not therefore do otherwise than
        arouse two servant-maids and go. She proceeded, with Pao-ch'ai,}
}
{
}

\Verse{136}
{
    \zA{宝钗笑道：“就是为那话了。”}
}
{
    \eA{out of the I Hung court, and then repaired all alone to lady Feng's on this
        side.}
}
{
}

\Verse{137}
{
    \zA{袭 人只得叫起两个丫头来，同着宝钗出怡红院，自往凤姐这里来。}
}
{
    \eA{It was indeed to communicate to her what had been decided about her, and to
        explain to her, as well, that though she could go and prostrate herself
        before Madame Wang, she could dispense with seeing dowager lady Chia.}
}
{
}

\Verse{138}
{
    \zA{果然是告诉他这话， 又教他给王夫人磕头，且不必去见贾母。}
}
{
    \eA{This news made Hsi Jen feel very awkward; to such an extent, that no sooner
        had she got through her visit to Madame Wang, than she returned in a hurry
        to her rooms. Pao-yü had already awoke.}
}
{
}

\Verse{139}
{
    \zA{倒把袭人说的甚觉不好意思。}
}
{
    \eA{He asked the reason why she had been called away, but Hsi Jen temporised by
        giving him an evasive answer.}
}
{
}

\Verse{140}
{
    \zA{及见过王夫人回来，宝玉已醒，问起原故，袭人且含糊答应。}
}
{
    \eA{And only at night, when every one was quiet, did Hsi Jen at length give him
        a full account of the whole matter. Pao-yü was delighted beyond measure.
        "I'll see now," he said,}
}
{
}

\Verse{141}
{
    \zA{至夜间人静，袭 人方告诉了。}
}
{
    \eA{with a face beaming with smiles, "whether you'll go back home or not.}
}
{
}

\Verse{142}
{
    \zA{宝玉喜不自禁，又向他笑道：“我可看你回家去不去了!那一回往家 里走了一趟，回来就说你哥哥要赎你，又说在这里没着落，终久算什么，说那些无
        情无义的生分话唬我。}
}
{
    \eA{On your return, after your last visit to your people, you stated that your
        brother wished to redeem you, adding that this place was no home for you,
        and that you didn't know what would become of you in the long run. You
        freely uttered all that language devoid of feeling and reason, and enough
        too to produce an estrangement between us,}
}
{
}

\Verse{143}
{
    \zA{从今我可看谁来敢叫你去？”}
}
{
    \eA{in order to frighten me; but I'd like to see who'll henceforward have the
        audacity to come and ask you to leave!"}
}
{
}

\Verse{144}
{
    \zA{袭人听了，冷笑道：“你倒别 这么说。}
}
{
    \eA{Hsi Jen, upon hearing this, smiled a smile full of irony. "You shouldn't say
        such things!"}
}
{
}

\Verse{145}
{
    \zA{从此以后，我是太太的人了，我要走，连你也不必告诉，只回了太太就走。”}
}
{
    \eA{she replied. "From henceforward I shall be our Madame Wang's servant, so
        that, if I choose to go I needn't even breathe a word to you. All I'll have
        to do will be to tell her, and then I shall be free to do as I like."}
}
{
}

\Verse{146}
{
    \zA{宝玉笑道：“就算我不好，你回了太太去了，叫别人听见说我不好，你去了，你有 什么意思呢？”}
}
{
    \eA{"But supposing that I behaved improperly," demurred Pao-yü laughingly, "and
        that you took your leave after letting mother know, you yourself will be
        placed in no nice fix, when people get wind that you left on account of my
        having been improper."}
}
{
}

\Verse{147}
{
    \zA{袭人笑道：“有什么没意思的？}
}
{
    \eA{"What no nice fix!" smiled Hsi Jen. "Is it likely that I am bound to serve
        even highway robbers?}
}
{
}

\Verse{148}
{
    \zA{难道下流人我也跟着罢？}
}
{
    \eA{Well, failing anything else, I can die; for human beings may live a hundred
        years,}
}
{
}

\Verse{149}
{
    \zA{再不然还有 个死呢!人活百岁，横竖要死，这口气没了，听不见看不见就罢了。”}
}
{
    \eA{but they're bound, in the long run, to fall a victim to death! And when this
        breath shall have departed, and I shall have lost the sense of hearing and
        of seeing, all will then be well!" When her rejoinder fell on his ear,}
}
{
}

\Verse{150}
{
    \zA{宝玉听见这 话，便忙握他的嘴，说道：“罢罢，你别说这些话了。”}
}
{
    \eA{Pao-yü promptly stopped her mouth with both his hands. "Enough! enough! that
        will do," he shouted. "There's no necessity for you to utter language of
        this kind."}
}
{
}

\Verse{151}
{
    \zA{袭人深知宝玉性情古怪， 听见奉承吉利话，又厌虚而不实，听了这些近情的实话，又生悲感。}
}
{
    \eA{Hsi Jen was well aware that Pao-yü was gifted with such a peculiar
        temperament, that he even looked upon flattering or auspicious phrases with
        utter aversion, treating them as meaningless and consequently insincere, so
        when, after listening to those truths, she had spoken with such pathos,}
}
{
}

\Verse{152}
{
    \zA{也后悔自己冒 撞，连忙笑着，用话截开，只拣宝玉那素日喜欢的，说些春风秋月，粉淡脂红，然 后又说到女儿如何好。}
}
{
    \eA{he, lapsed into another of his melancholy moods, she blamed herself for the
        want of consideration she had betrayed. Hastily therefore putting on a
        smile, she tried to hit upon some suitable remarks, with which to interrupt
        the conversation. Her choice fell upon those licentious and immodest topics,}
}
{
}

\Verse{153}
{
    \zA{不觉又说到女儿死的上头，袭人忙掩住口。}
}
{
    \eA{which had ever been a relish to the taste of Pao-yü; and from these the
        conversation drifted to the subject of womankind. But when, subsequently,
        reference was made to the excellency of the weak sex, they somehow or other
        also came to touch upon the mortal nature of women, and Hsi Jen promptly
        closed her lips in silence.}
}
{
}

\Verse{154}
{
    \zA{宝玉听至浓快处，见他不说了，便笑道：“人谁不死？}
}
{
    \eA{Noticing however that now that the conversation had reached a point so full
        of zest for him, she had nothing to say for herself, Pao-yü smilingly
        remarked:}
}
{
}

\Verse{155}
{
    \zA{只要死的好。}
}
{
    \eA{"What human being is there that can escape death?}
}
{
}

\Verse{156}
{
    \zA{那些须眉 浊物只听见‘文死谏’‘武死战’这二死是大丈夫的名节，便只管胡闹起来。}
}
{
    \eA{But the main thing is to come to a proper end! All that those abject male
        creatures excel in is, the civil officers, to sacrifice their lives by
        remonstrating with the Emperor; and, the military, to leave their bones on
        the battlefield.}
}
{
}

\Verse{157}
{
    \zA{那里 知道有昏君，方有死谏之臣，只顾他邀名，猛拚一死，将来置君父于何地？}
}
{
    \eA{Both these deaths do confer, after life is extinct, the fame of great men
        upon them; but isn't it, in fact, better for them not to die? For as it is
        absolutely necessary that there should be a disorderly Emperor before they
        can afford any admonition,}
}
{
}

\Verse{158}
{
    \zA{必定有 刀兵，方有死战，他只顾图汗马之功，猛拚一死，将来弃国于何地？”}
}
{
    \eA{to what future fate do they thus expose their sovereign, if they rashly
        throw away their lives, with the sole aim of reaping a fair name for
        themselves? War too must supervene before they can fight; but if they go and
        recklessly lay down their lives,}
}
{
}

\Verse{159}
{
    \zA{袭人不等说 完，便道：“古时候儿这些人，也因出于不得已他才死啊。”}
}
{
    \eA{with the exclusive idea of gaining the reputation of intrepid warriors, to
        what destiny will they abandon their country by and bye? Hence it is that
        neither of these deaths can be looked upon as a legitimate death." "Loyal
        ministers," Hsi Jen argued,}
}
{
}

\Verse{160}
{
    \zA{宝玉道：“那武将要 是疏谋少略的，他自己无能，白送了性命，这难道也是不得已么？}
}
{
    \eA{"and excellent generals simply die because it isn't in their power to do
        otherwise." "Military officers," Pao-yü explained, "place such entire
        reliance upon brute force that they become lax in their stratagems and
        faulty in their plans.}
}
{
}

\Verse{161}
{
    \zA{那文官更不比武 官了：他念两句书，记在心里，若朝廷少有瑕疵，他就胡弹乱谏，邀忠烈之名；倘 有不合，浊气一涌，即时拚死，这难道也是不得已？}
}
{
    \eA{It's because they don't possess any inherent abilities that they lose their
        lives. Could one therefore, pray, say that they had no other alternative?
        Civil officials, on the other hand, can still less compare with military
        officers. They read a few passages from books,}
}
{
}

\Verse{162}
{
    \zA{要知道那朝廷是受命于天，若 非圣人，那天也断断不把这万几重任交代。}
}
{
    \eA{and commit them to memory; and, on the slightest mistake made by the
        Emperor, they're at once rash enough to remonstrate with him, prompted by
        the sole idea of attaining the fame of loyalty and devotion.}
}
{
}

\Verse{163}
{
    \zA{可知那些死的，都是沽名钓誉，并不知 君臣的大义。}
}
{
    \eA{But, as soon as their stupid notions have bubbled over, they forfeit their
        lives, and is it likely that it doesn't lie within their power to do
        otherwise?}
}
{
}

\Verse{164}
{
    \zA{比如我此时若果有造化，趁着你们都在眼前，我就死了，再能够你们 哭我的眼泪，流成大河，把我的尸首漂起来，送到那鸦雀不到的幽僻去处，随风化
        了，自此再不托生为人，这就是我死的得时了。”}
}
{
    \eA{Why, they should also bear in mind that the Emperor receives his decrees
        from Heaven; and, that were he not a perfect man, Heaven itself would, on no
        account whatever, confer upon him a charge so extremely onerous. This makes
        it evident therefore that the whole pack and parcel of those officers, who
        are dead and gone, have invariably fallen victims to their endeavours to
        attain a high reputation,}
}
{
}

\Verse{165}
{
    \zA{袭人忽见说出这些疯话来，忙说： “困了。”}
}
{
    \eA{and that they had no knowledge whatever of the import of the great principle
        of right!}
}
{
}

\Verse{166}
{
    \zA{不再答言。}
}
{
    \eA{Take me as an instance now.}
}
{
}

\Verse{167}
{
    \zA{那宝玉方合眼睡着。}
}
{
    \eA{Were really mine the good fortune of departing life at a fit time,}
}
{
}

\Verse{168}
{
    \zA{次日也就丢开。}
}
{
    \eA{I'd avail myself of the present when all you girls are alive,}
}
{
}

\Verse{169}
{
    \zA{一日，宝玉因各处游的腻烦，便想起《牡丹亭》曲子来，自己看了两遍，犹不 惬怀，因闻得梨香院的十二个女孩儿中，有个小旦龄官，唱的最妙。}
}
{
    \eA{to pass away. And could I get you to shed such profuse tears for me as to
        swell out into a stream large enough to raise my corpse and carry it to some
        secluded place, whither no bird even has ever wended its flight, and could I
        become invisible like the wind, and nevermore from this time,}
}
{
}

\Verse{170}
{
    \zA{因出了角门来 找时，只见葵官药官都在院内，见宝玉来了，都笑迎让坐。}
}
{
    \eA{come into existence as a human being, I shall then have died at a proper
        season." Hsi Jen suddenly awoke to the fact that he was beginning to give
        vent to a lot of twaddle, and speedily, pleading fatigue,}
}
{
}

\Verse{171}
{
    \zA{宝玉因问：“龄官在那 里？”}
}
{
    \eA{she paid no further notice to him. This compelled Pao-yü to at last be quiet
        and go to sleep.}
}
{
}

\Verse{172}
{
    \zA{都告诉他说：“在他屋里呢。”}
}
{
    \eA{By the morrow, all recollection of the discussion had vanished from his
        mind.}
}
{
}

\Verse{173}
{
    \zA{宝玉忙至他屋内，只见龄官独自躺在枕上， 见他进来，动也不动。}
}
{
    \eA{One day, Pao-yü was feeling weary at heart, after strolling all over the
        place, when remembering the song of the "Peony Pavilion," he read it over
        twice to himself;}
}
{
}

\Verse{174}
{
    \zA{宝玉身旁坐下，因素昔与别的女孩子玩惯了的，只当龄官也 和别人一样，遂近前陪笑，央他起来唱一套“袅晴丝”。}
}
{
    \eA{but still his spirits continued anything but joyous. Having heard, however,
        that among the twelve girls in the Pear Fragrance Court there was one called
        Ling Kuan, who excelled in singing, he purposely issued forth by a side gate
        and came in search of her. But the moment he got there,}
}
{
}

\Verse{175}
{
    \zA{不想龄官见他坐下，忙抬 起身来躲避，正色说道：“嗓子哑了，前儿娘娘传进我们去，我还没有唱呢。”}
}
{
    \eA{he discovered Pao Kuan, and Yü Kuan in the court. As soon as they caught
        sight of Pao-yü, they, with one consent, smiled and urged him to take a
        seat. Pao-yü then inquired where Ling Kuan was. Both girls explained that
        she was in her room,}
}
{
}

\Verse{176}
{
    \zA{宝 玉见他坐正了，再一细看，原来就是那日蔷薇花下画“蔷”字的那一个。}
}
{
    \eA{so Pao-yü hastened in. Here he found Ling Kuan alone, reclining against a
        pillow. Though perfectly conscious of his arrival, she did not move a
        muscle. Pao-yü ensconced himself next to her.}
}
{
}

\Verse{177}
{
    \zA{又见如此 景况，从来未经过这样被人弃厌，自己便讪讪的，红了脸，只得出来了。}
}
{
    \eA{He had always been in the habit of playing with the rest of the girls, so
        thinking that Ling Kuan was like the others, he felt impelled to draw near
        her and to entreat her, with a forced smile, to get up and sing part of the
        "Niao Ch'ing Ssu."}
}
{
}

\Verse{178}
{
    \zA{药官等不解何故，因问其所以，宝玉便告诉了他。}
}
{
    \eA{But his hopes were baffled; for as soon as Ling Kuan perceived him sit down,
        she impetuously raised herself and withdrew from his side.}
}
{
}

\Verse{179}
{
    \zA{宝官笑说道：“只略等一等， 蔷二爷来了，他叫唱是必唱的。”}
}
{
    \eA{"I'm hoarse," she rejoined with a stern expression on her face. "The Empress
        the other day called us into the palace; but I couldn't sing even then."
        Seeing her sit bolt upright,}
}
{
}

\Verse{180}
{
    \zA{宝玉听了，心下纳闷，因问：“蔷哥儿那里去了？”}
}
{
    \eA{Pao-yü went on to pass her under a minute survey. He discovered that it was
        the girl, whom he had, some time ago beheld under the cinnamon roses,}
}
{
}

\Verse{181}
{
    \zA{宝官道：“才出去了，一定就是龄官儿要什么，他去变弄去了。”}
}
{
    \eA{drawing the character "Ch'iang." But seeing the reception she accorded him,
        who had never so far known what it was to be treated contemptuously by any
        one,}
}
{
}

\Verse{182}
{
    \zA{宝玉听了以为奇 特。}
}
{
    \eA{he blushed crimson, while muttering some abuse to himself,}
}
{
}

\Verse{183}
{
    \zA{少站片时，果见贾蔷从外头来了，手里提着个雀儿笼子，上面扎着小戏台，并 一个雀儿，兴兴头头往里来找龄官。}
}
{
    \eA{and felt constrained to quit the room. Pao Kuan and her companion could not
        fathom why he was so red and inquired of him the reason. Pao-yü told them.
        "Wait a while," Pao Kuan said, "until Mr. Ch'iang Secundus comes; and when
        he asks her to sing, she is bound to sing." Pao-yü at these words felt very
        sad within himself.}
}
{
}

\Verse{184}
{
    \zA{见了宝玉，只得站住。}
}
{
    \eA{"Where's brother Ch'iang gone to?" he asked. "He's just gone out,"}
}
{
}

\Verse{185}
{
    \zA{宝玉问他：“是个什么 雀儿？”}
}
{
    \eA{Pao Kuan answered. "Of course, Ling Kuan must have wanted something or
        other,}
}
{
}

\Verse{186}
{
    \zA{贾蔷笑道：“是个玉顶儿，还会衔旗串戏。”}
}
{
    \eA{and he's gone to devise ways and means to bring it to her." Pao-yü thought
        this remark very extraordinary.}
}
{
}

\Verse{187}
{
    \zA{宝玉道：“多少钱买的？”}
}
{
    \eA{But after standing about for a while, he actually saw Chia Ch'iang arrive
        from outside,}
}
{
}

\Verse{188}
{
    \zA{贾蔷道：“一两八钱银子。”}
}
{
    \eA{carrying a cage, with a tiny stage inserted at the top, and a bird as well;
        and wend his steps,}
}
{
}

\Verse{189}
{
    \zA{一面说，一面让宝玉坐，自己往龄官屋里来。}
}
{
    \eA{in a gleeful mood, towards the interior to join Ling Kuan. The moment,
        however, he noticed Pao-yü, he felt under the necessity of halting.}
}
{
}

\Verse{190}
{
    \zA{宝玉此刻把听曲子的心都没了，且要看他和龄官是怎么样。}
}
{
    \eA{"What kind of bird is that?" Pao-yü asked. "Can it hold a flag in its beak,
        or do any tricks?" "It's the 'jade-crested and gold-headed bird,'" smiled
        Chia Ch'iang. "How much did you give for it?"}
}
{
}

\Verse{191}
{
    \zA{只见贾蔷进去，笑 道：“你来瞧这个玩意儿。”}
}
{
    \eA{Pao-yü continued. "A tael and eight mace," replied Chia Ch'iang. But while
        replying to his inquiries, he motioned to Pao-yü to take a seat,}
}
{
}

\Verse{192}
{
    \zA{龄官起身问：“是什么？”}
}
{
    \eA{and then went himself into Ling Kuan's apartment. Pao-yü had, by this time,}
}
{
}

\Verse{193}
{
    \zA{贾蔷道：“买了个雀儿给 你玩，省了你天天儿发闷。}
}
{
    \eA{lost every wish of hearing a song. His sole desire was to find what
        relations existed between his cousin and Ling Kuan, when he perceived Chia
        Ch'iang walk in and laughingly say to her,}
}
{
}

\Verse{194}
{
    \zA{我先玩个你瞧瞧。”}
}
{
    \eA{"Come and see this thing." "What's it?" Ling Kuan asked,}
}
{
}

\Verse{195}
{
    \zA{说着，便拿些谷子，哄的那个雀儿 果然在那戏台上衔着鬼脸儿和旗帜乱串。}
}
{
    \eA{rising. "I've bought a bird for you to amuse yourself with," Chia Ch'iang
        added, "so that you mayn't daily feel dull and have nothing to distract
        yourself with. But I'll first play with it and let you see." With this
        prelude, he took a few seeds and began to coax the bird, until it, in point
        of fact,}
}
{
}

\Verse{196}
{
    \zA{众女孩子都笑了，独龄官冷笑两声，赌气 仍睡着去了。}
}
{
    \eA{performed various tricks, on the stage, clasping in its beak a mask and a
        flag. All the girls shouted out: "How nice;" with the sole exception of Ling
        Kuan,}
}
{
}

\Verse{197}
{
    \zA{贾蔷还只管陪笑问他：“好不好？”}
}
{
    \eA{who gave a couple of apathetic smirks, and went in a huff to lie down. Again
        Chia Ch'iang, however,}
}
{
}

\Verse{198}
{
    \zA{龄官道：“你们家把好好儿的人 弄了来，关在这牢坑里，学这个还不算，你这会子又弄个雀儿来，也干这个浪事! 你分明弄了来打趣形容我们，还问‘好不好’！”}
}
{
    \eA{kept on forcing smiles, and inquiring of her whether she liked it or not.
        "Isn't it enough," Ling Kuan observed, "that your family entraps a fine lot
        of human beings like us and coops us up in this hole to study this stuff and
        nonsense, but do you also now go and get a bird, which likewise is, as it
        happens,}
}
{
}

\Verse{199}
{
    \zA{贾蔷听了，不觉站起来，连忙赌 神起誓，又道：“今儿我那里的糊涂油蒙了心，费一二两银子买他，原说解闷儿， 就没想到这上头。}
}
{
    \eA{up to this sort of thing? You distinctly fetch it to make fun of us, and
        mimick us, and do you still ask me whether I like it or not?" Hearing this
        reproach, Chia Ch'iang of a sudden sprang to his feet with alacrity and
        vehemently endeavoured by vowing and swearing to establish his innocence.
        "How ever could I have been such a fool to-day,"}
}
{
}

\Verse{200}
{
    \zA{罢了，放了生，倒也免你的灾。”}
}
{
    \eA{he proceeded, "as to go and throw away a tael or two to purchase this bird?
        I really did it in the hope that it would afford you amusement.}
}
{
}

\Verse{201}
{
    \zA{说着，果然将那雀儿放了，一 顿把那笼子拆了。}
}
{
    \eA{I never for a moment entertained such thoughts as those you credit me with.
        But never mind; I'll let it go, and save you all this misery!"}
}
{
}

\Verse{202}
{
    \zA{龄官还说：“那雀儿虽不如人，他也有个老雀儿在窝里，你拿了 他来，弄这个劳什子，也忍得？}
}
{
    \eA{So saying, he verily gave the bird its liberty; and, with one blow, he
        smashed the cage to atoms. "This bird," still argued Ling Kuan, "differs,
        it's true, from a human being; but it too has a mother and father in its
        nest,}
}
{
}

\Verse{203}
{
    \zA{今儿我咳嗽出两口血来，太太打发人来找你，叫你 请大夫来细问问，你且弄这个来取笑儿。}
}
{
    \eA{and could you have had the heart to bring it here to perform these silly
        pranks? In coughing to-day, I expectorated two mouthfuls of blood, and
        Madame Wang sent some one here to find you so as to tell you to ask the
        doctor round to minutely diagnose my complaint,}
}
{
}

\Verse{204}
{
    \zA{偏是我这没人管没人理的，又偏爱害病！”}
}
{
    \eA{and have you instead brought this to mock me with? But it so happens that I,
        who have not a soul to look after me,}
}
{
}

\Verse{205}
{
    \zA{贾蔷听说，连忙说道：“昨儿晚上我问了大夫，他说，‘不相干，吃两剂药，后儿 再瞧。’}
}
{
    \eA{or to care for me, also have the fate to fall ill!" Chia Ch'iang listened to
        her. "Yesterday evening," he eagerly explained, "I asked the doctor about
        it. He said that it was nothing at all, that you should take a few doses of
        medicine,}
}
{
}

\Verse{206}
{
    \zA{谁知今儿又吐了？}
}
{
    \eA{and that he would be coming again in a day or two to see how you were
        getting on.}
}
{
}

\Verse{207}
{
    \zA{这会子就请他去。”}
}
{
    \eA{But who'd have thought it, you have again to-day expectorated blood.}
}
{
}

\Verse{208}
{
    \zA{说着便要请去。}
}
{
    \eA{I'll go at once and invite him to come round."}
}
{
}

\Verse{209}
{
    \zA{龄官又叫：“站住， 这会子大毒日头地下，你赌气去请了来，我也不瞧。”}
}
{
    \eA{Speaking the while, he was about to go immediately when Ling Kuan cried out
        and stopped him. "Do you go off in a tantrum in this hot broiling sun?" she
        said. "You may ask him to come,}
}
{
}

\Verse{210}
{
    \zA{贾蔷听如此说，只得又站住。}
}
{
    \eA{but I won't see him." When he heard her resolution, Chia Ch'iang had
        perforce to stand still.}
}
{
}

\Verse{211}
{
    \zA{宝玉见了这般景况，不觉痴了。}
}
{
    \eA{Pao-yü, perceiving what transpired between them, fell unwittingly in a dull
        reverie.}
}
{
}

\Verse{212}
{
    \zA{这才领会过画“蔷”深意。}
}
{
    \eA{He then at length got an insight into the deep import of the tracing of the
        character "Ch'iang."}
}
{
}

\Verse{213}
{
    \zA{自己站不住，便抽 身走了。}
}
{
    \eA{But unable to bear the ordeal any longer, he forthwith took himself out of
        the way.}
}
{
}

\Verse{214}
{
    \zA{贾蔷一心都在龄官身上，竟不曾理会，倒是别的女孩子送出来了。}
}
{
    \eA{So absorbed, however, was Chia Ch'iang's whole mind with Ling Kuan that he
        could not even give a thought to escorting any one; and it was, in fact, the
        rest of the singing-girls who saw (Pao-yü) out.}
}
{
}

\Verse{215}
{
    \zA{那宝玉 一心裁夺盘算，痴痴的回至怡红院中，正值黛玉和袭人坐着说话儿呢。}
}
{
    \eA{Pao-yü's heart was gnawed with doubts and conjectures. In an imbecile frame
        of mind, he came to the I Hung court. Lin Tai-yü was, at the moment, sitting
        with Hsi Jen, and chatting with her.}
}
{
}

\Verse{216}
{
    \zA{宝玉一进来， 就和袭人长叹，说道：“我昨儿晚上的话，竟说错了，怪不得老爷说我是‘管窥蠡 测’!昨夜说你们的眼泪单葬我，这就错了。}
}
{
    \eA{As soon as Pao-yü entered his quarters, he addressed himself to Hsi Jen,
        with a long sigh. "I was very wrong in what I said yesterday evening," he
        remarked. "It's no matter of surprise that father says that I am so
        narrow-minded that I look at things through a tube and measure them with a
        clam-shell. I mentioned something last night about having nothing but tears,}
}
{
}

\Verse{217}
{
    \zA{看来我竟不能全得。}
}
{
    \eA{shed by all of you girls, to be buried in. But this was a mere delusion!}
}
{
}

\Verse{218}
{
    \zA{从此后，只好各 人得各人的眼泪罢了。”}
}
{
    \eA{So as I can't get the tears of the whole lot of you, each one of you can
        henceforward keep her own for herself,}
}
{
}

\Verse{219}
{
    \zA{袭人只道昨夜不过是些玩话，已经忘了，不想宝玉又提起 来，便笑道：“你可真真有些个疯了！”}
}
{
    \eA{and have done." Hsi Jen had flattered herself that the words he had uttered
        the previous evening amounted to idle talk, and she had long ago dispelled
        all thought of them from her mind, but when Pao-yü unawares made further
        allusion to them, she smilingly rejoined: "You are verily somewhat cracked!"}
}
{
}

\Verse{220}
{
    \zA{宝玉默默不对。}
}
{
    \eA{Pao-yü kept silent, and attempted to make no reply.}
}
{
}

\Verse{221}
{
    \zA{自此深悟人生情缘，各有 分定，只是每每暗伤：“不知将来葬我洒泪者为谁？”}
}
{
    \eA{Yet from this time he fully apprehended that the lot of human affections is,
        in every instance, subject to predestination, and time and again he was wont
        to secretly muse, with much anguish: "Who,}
}
{
}

\Verse{222}
{
    \zA{且说黛玉当下见宝玉如此形象，便知是又从那里着了魔来，也不便多问，因说 道：“我才在舅母跟前，听见说明儿是薛姨妈的生日，叫我顺便来问你出去不出去。}
}
{
    \eA{I wonder, will shed tears for me, at my burial?" Lin Tai-yü, for we will now
        allude to her, noticed Pao-yü's behaviour, but readily concluding that he
        must have been, somewhere or other, once more possessed by some malignant
        spirit,}
}
{
}

\Verse{223}
{
    \zA{你打发人前头说一声去。”}
}
{
    \eA{she did not feel it advisable to ask many questions. "I just saw," she
        consequently observed,}
}
{
}

\Verse{224}
{
    \zA{宝玉道：“上回连大老爷的生日我也没去，这会子我又 去，倘或碰见了人呢？}
}
{
    \eA{"my maternal aunt, who hearing that to-morrow is Miss Hsüeh's birthday, bade
        me come at my convenience to ask you whether you'll go or not, (and to tell
        you) to send some one ahead to let them know what you mean to do."}
}
{
}

\Verse{225}
{
    \zA{我一概都不去。}
}
{
    \eA{"I didn't go the other day, when it was Mr. Chia She's birthday,}
}
{
}

\Verse{226}
{
    \zA{这么怪热的，又穿衣裳!我不去，姨妈也未必 恼。”}
}
{
    \eA{so I won't go now." Pao-yü answered. "If it is a matter of meeting any one,
        I won't go anywhere. On a hot day like this to again don my ceremonial
        dress!}
}
{
}

\Verse{227}
{
    \zA{袭人忙道：“这是什么话？}
}
{
    \eA{No, I won't go. Aunt is not likely to feel displeased with me!" "What are
        you driving at?" Hsi Jen speedily ventured.}
}
{
}

\Verse{228}
{
    \zA{他比不得大老爷。}
}
{
    \eA{"She couldn't be put on the same footing as our senior master!}
}
{
}

\Verse{229}
{
    \zA{这里又住的近，又是亲戚，你 不去，岂不叫他思量？}
}
{
    \eA{She lives close by here. Besides she's a relative. Why, if you don't go,
        won't you make her imagine things?}
}
{
}

\Verse{230}
{
    \zA{你怕热，就清早起来，到那里磕个头、吃钟茶再来，岂不好 看？”}
}
{
    \eA{Well, if you dread the heat, just get up at an early hour and go over and
        prostrate yourself before her, and come back again, after you've had a cup
        of tea.}
}
{
}

\Verse{231}
{
    \zA{宝玉尚未说话，黛玉便先笑道：“你看着人家赶蚊子的分上，也该去走走。”}
}
{
    \eA{Won't this look well?" Before Pao-yü had time to say anything by way of
        response, Tai-yü anticipated him. "You should," she smiled, "go as far as
        there for the sake of her, who drives the mosquitoes away from you."}
}
{
}

\Verse{232}
{
    \zA{宝玉不解，忙问：“怎么赶蚊子？”}
}
{
    \eA{Pao-yü could not make out the drift of her insinuation. "What about driving
        mosquitoes away?"}
}
{
}

\Verse{233}
{
    \zA{袭人便将昨日睡觉无人作伴，宝姑娘坐了一坐 的话，告诉宝玉。}
}
{
    \eA{he vehemently inquired. Hsi Jen then explained to him how while he was fast
        asleep the previous day and no one was about to keep him company, Miss
        Pao-ch'ai had sat with him for a while. "It shouldn't have been done!"
        Pao-yü promptly exclaimed,}
}
{
}

\Verse{234}
{
    \zA{宝玉听了，忙说：“不该!我怎么睡着了？}
}
{
    \eA{after hearing her explanations. "But how did I manage to go to sleep and
        show such utter discourtesy to her? I must go to-morrow!" he then went on to
        add.}
}
{
}

\Verse{235}
{
    \zA{就亵渎了他！”}
}
{
    \eA{But while these words were still on his lips,}
}
{
}

\Verse{236}
{
    \zA{一面又 说：“明日必去。”}
}
{
    \eA{he unexpectedly caught sight of Shih Hsian-yün walk in in full dress,}
}
{
}

\Verse{237}
{
    \zA{正说着，忽见湘云穿得齐齐整整的走来，辞说家里打发人来接他。}
}
{
    \eA{to bid them adieu, as she said that some one had been sent from her home to
        fetch her away. The moment Pao-yü and Tai-yü heard what was the object of
        her visit, they quickly rose to their feet and pressed her to take a seat.}
}
{
}

\Verse{238}
{
    \zA{宝玉黛玉听 说，忙站起来让坐，湘云也不坐，宝黛两个只得送他至前面。}
}
{
    \eA{But Shih Hsiang-yün would not sit down, so Pao-yü and Tai-yü were compelled
        to escort her as far as the front part of the mansion. Shih Hsiang-yün's
        eyes were brimming with tears;}
}
{
}

\Verse{239}
{
    \zA{那湘云只是眼泪汪汪 的，见有他家的人在跟前，又不敢十分委屈。}
}
{
    \eA{but realising that several people from her home were present, she did not
        have the courage to give full vent to her feelings. But when shortly
        Pao-ch'ai ran over to find her,}
}
{
}

\Verse{240}
{
    \zA{少时宝钗赶来，愈觉缱绻难舍。}
}
{
    \eA{she felt so much the more drawn towards them, that she could not brook to
        part from them.}
}
{
}

\Verse{241}
{
    \zA{还是 宝钗心内明白，他家里人若回去告诉了他婶娘，待他家去了，又恐怕他受气，因此 倒催着他走了。}
}
{
    \eA{Pao-ch'ai, however, inwardly understood that if her people told her aunt
        anything on their return, there would again be every fear of her being blown
        up, as soon as she got back home, and she therefore urged her to start on
        her way. One and all then walked with her up to the second gate,}
}
{
}

\Verse{242}
{
    \zA{众人送至二门前，宝玉还要往外送他，倒是湘云拦住了。}
}
{
    \eA{and Pao-yü wished to accompany her still further outside, but Shih
        Hsiang-yün deterred him. Presently, they turned to go back. But once more,
        she called Pao-yü to her,}
}
{
}

\Verse{243}
{
    \zA{一时，回 身又叫宝玉到跟前，悄悄的嘱咐道：“就是老太太想不起我来，你时常提着，好等 老太太打发人接我去。”}
}
{
    \eA{and whispered to him in a soft tone of voice: "Should our venerable senior
        not think of me do often allude to me, so that she should depute some one to
        fetch me." Pao-yü time after time assured her that he would comply with her
        wishes.}
}
{
}

\Verse{244}
{
    \zA{宝玉连连答应了。}
}
{
    \eA{And having followed her with their eyes,}
}
{
}

\Verse{245}
{
    \zA{眼看着他上车去了，大家方才进来。}
}
{
    \eA{while she got into her curricle and started, they eventually retraced their
        steps towards the inner compound.}
}
{
}

\Verse{246}
{
    \zA{要知端底，且看下回分解。}
}
{
    \eA{But, reader, if you like to follow up the story,}
}
{
}

\Verse{247}
{
    \zA{(本章完)}
}
{
    \eA{peruse the details contained in the chapter below.}
}
{
}
